\chapter{Herleitung der Methode der kleinsten Quadrate}\label{sec:appendix_methode_der_kleinsten_quadrate}

\section*{Definition der Zielfunktion}
Das Ziel der linearen Regression besteht darin, die bestmögliche Gerade $y = mx + q$ zu finden, die die Summe der quadrierten Residuen zwischen den vorhergesagten Werten und den tatsächlichen Datenpunkten minimiert. Die Summe der quadrierten Abstände zwischen den Datenpunkten und der Geraden wird definiert als:
$$
S = \sum_{i=1}^n (y_i - (mx_i + q))^2
$$

\section*{Partielle Ableitungen}
Um $S$ zu minimieren, nehmen wir die partiellen Ableitungen bezüglich $m$ und $q$ und setzen diese gleich Null.

\subsection*{Partielle Ableitung nach $m$}
$$
\frac{\partial S}{\partial m} = \frac{\partial}{\partial m} \sum_{i=1}^n (y_i - mx_i - q)^2
$$
$$
= \sum_{i=1}^n 2(y_i - mx_i - q)(-x_i)
$$
$$
= -2 \sum_{i=1}^n x_i(y_i - mx_i - q)
$$
Auf Null setzen für die Minimierung:
$$
\sum x_iy_i - m\sum x_i^2 - q\sum x_i = 0
$$
$$
\sum x_iy_i = m\sum x_i^2 + q\sum x_i \quad \text{(Gleichung 1)}
$$

\subsection*{Partielle Ableitung nach $q$}
$$
\frac{\partial S}{\partial q} = \frac{\partial}{\partial q} \sum_{i=1}^n (y_i - mx_i - q)^2
$$
$$
= \sum_{i=1}^n 2(y_i - mx_i - q)(-1)
$$
$$
= -2 \sum_{i=1}^n (y_i - mx_i - q)
$$
Auf Null setzen für die Minimierung:
$$
\sum y_i - m\sum x_i - nq = 0
$$
$$
\sum y_i = m\sum x_i + nq \quad \text{(Gleichung 2)}
$$

\section*{Lösung der Gleichungen}
Aus Gleichung 2:
$$
q = \frac{\sum y_i - m\sum x_i}{n} \quad \text{(Gleichung 3)}
$$
Einsetzen der Gleichung 3 in Gleichung 1:
$$
\sum x_iy_i = m\sum x_i^2 + \left(\frac{\sum y_i - m\sum x_i}{n}\right)\sum x_i
$$
$$
\sum x_iy_i = m\sum x_i^2 + \frac{(\sum y_i)\sum x_i}{n} - \frac{m(\sum x_i)^2}{n}
$$
Auflösen nach $m$:
$$
m = \frac{n\sum x_iy_i - \sum x_i \sum y_i}{n\sum x_i^2 - (\sum x_i)^2}
$$
Bestimmen von $q$ über Einsetzen von $m$ in Gleichung 3:
$$
q = \frac{\sum y_i - m\sum x_i}{n}
$$


\chapter{Formelle Definition von Markov-Ketten}\label{ch:markov-formal}

Formal definieren wir eine Markov-Kette über einen diskreten Zustandsraum $S$ und eine Übergangsmatrix $P$.

\begin{mydefinition}[Markov-Kette]
Eine \textbf{Markov-Kette} ist ein stochastischer Prozess $\{X_n : n \in \mathbb{N}_0\}$ mit Werten in einem abzählbaren Zustandsraum $S$, der die Markov-Eigenschaft erfüllt:
\[
    P(X_{n+1} = j | X_n = i, X_{n-1} = i_{n-1}, \ldots, X_0 = i_0) = P(X_{n+1} = j | X_n = i)
\]

Die Wahrscheinlichkeit $p_{ij} = P(X_{n+1} = j | X_n = i)$ bezeichnet die Übergangswahrscheinlichkeit von Zustand $i$ zu Zustand $j$. Die Matrix $P = (p_{ij})_{i,j \in S}$ heisst \textbf{Übergangsmatrix} der Markov-Kette.

Eine Markov-Kette heisst \textbf{homogen}, wenn die Übergangswahrscheinlichkeiten nicht vom Zeitpunkt $n$ abhängen, sondern nur von den beteiligten Zuständen $i$ und $j$.
\end{mydefinition}

Der Zustandsvektor $\pi^{(n)} = (\pi_1^{(n)}, \pi_2^{(n)}, \ldots)$ gibt die Wahrscheinlichkeitsverteilung der Zustände zum Zeitpunkt $n$ an, wobei $\pi_i^{(n)} = P(X_n = i)$. Der Anfangszustand wird durch den Vektor $\pi^{(0)}$ beschrieben.

Die zeitliche Entwicklung einer Markov-Kette lässt sich durch wiederholte Multiplikation mit der Übergangsmatrix $P$ berechnen:
\[
    \pi^{(n+1)} = \pi^{(n)} \cdot P
\]

Das bedeutet, dass die Verteilung nach $n$ Schritten durch $\pi^{(n)} = \pi^{(0)} \cdot P^n$ gegeben ist.


\chapter{Eigenschaften von Markov-Ketten}\label{app:markov-properties}

Markov-Ketten und ihre Zustände lassen sich durch verschiedene Eigenschaften charakterisieren. Diese Eigenschaften helfen, das Langzeitverhalten der Kette besser zu verstehen.

\begin{itemize}
    \item \textbf{Aperiodizität:} Ein Zustand $i$ hat die Periode $k \geq 1$, wenn jede Rückkehr zu $i$ nur in einer Anzahl von Schritten möglich ist, die durch $k$ teilbar ist. Ist $k=1$, so heisst der Zustand \emph{aperiodisch}. Sind alle Zustände aperiodisch, so ist die Markov-Kette aperiodisch.
    \item \textbf{Irreduzibilität:} Eine Markov-Kette heisst \emph{irreduzibel}, wenn es zwischen jedem Zustandspaar $i$ und $j$ eine Kette von Schritten mit positiver Wahrscheinlichkeit gibt.
    \item \textbf{Absorbierende Zustände:} Ein Zustand $i$ ist \emph{absorbierend}, wenn $P_{i,i} = 1$ gilt, d.h. die Kette bleibt nach Erreichen dieses Zustands immer dort.
    \item \textbf{Rekurrenz und Transienz:} Ein Zustand ist \emph{rekurrent}, wenn die Kette mit Wahrscheinlichkeit $1$ irgendwann wieder in diesen Zustand zurückkehrt. Andernfalls ist er \emph{transient}. Ist die erwartete Rückkehrzeit endlich, heisst der Zustand \emph{positiv rekurrent}, sonst \emph{null-rekurrent}.
    \item \textbf{Ergodizität:} Ein Zustand ist \emph{ergodisch}, wenn er positiv rekurrent und aperiodisch ist. Eine Markov-Kette ist ergodisch, wenn alle ihre Zustände ergodisch sind.
\end{itemize}

\begin{mychallenge}[Aperiodizität]
    Gegeben sei die folgende Übergangsmatrix:
    \[
        P = \begin{pmatrix}
            0{,}5 & 0{,}5 \\
            0{,}5 & 0{,}5 \\
        \end{pmatrix}
    \]
    Untersuchen Sie, ob die Markov-Kette aperiodisch ist.

    \begin{myanswer}
        Von jedem Zustand kann man in einem Schritt in jeden anderen Zustand (auch in sich selbst) gelangen. Die Rückkehr zu einem Zustand ist also in jedem Schritt möglich, insbesondere in aufeinanderfolgenden Schritten. Daher ist die Periode $k=1$ und die Kette ist aperiodisch.
    \end{myanswer}
\end{mychallenge}

\begin{mychallenge}[Irreduzibilität]
    Gegeben sei die Übergangsmatrix:
    \[
        P = \begin{pmatrix}
            0 & 1 \\
            1 & 0 \\
        \end{pmatrix}
    \]
    Ist diese Markov-Kette irreduzibel? Ist sie aperiodisch?

    \begin{myanswer}
        Von jedem Zustand kann man in einem Schritt zum jeweils anderen Zustand gelangen, also ist die Kette irreduzibel. Die Rückkehr zu einem Zustand ist aber nur in einer geraden Anzahl von Schritten möglich (2, 4, 6, ...), daher ist die Periode $k=2$ und die Kette ist periodisch (nicht aperiodisch).
    \end{myanswer}
\end{mychallenge}

\begin{mychallenge}[Absorbierende Zustände]
    Gegeben sei die Übergangsmatrix:
    \[
        P = \begin{pmatrix}
            1     & 0     \\
            0{,}3 & 0{,}7 \\
        \end{pmatrix}
    \]
    Gibt es einen absorbierenden Zustand?

    \begin{myanswer}
        Der Zustand 1 ist absorbierend, da $P_{1,1} = 1$. Das heisst, wenn die Kette einmal in Zustand 1 ist, bleibt sie dort.
    \end{myanswer}
\end{mychallenge}

\begin{mychallenge}[Rekurrenz und Transienz]
    Betrachten Sie die folgende Übergangsmatrix:
    \[
        P = \begin{pmatrix}
            0{,}6 & 0{,}4 \\
            0{,}2 & 0{,}8 \\
        \end{pmatrix}
    \]
    Ist die Kette irreduzibel? Sind die Zustände rekurrent oder transient?

    \begin{myanswer}
        Von jedem Zustand kann man mit positiver Wahrscheinlichkeit in den anderen gelangen, also ist die Kette irreduzibel. Da es nur zwei Zustände gibt und man immer wieder zurückkehren kann, sind beide Zustände rekurrent.
    \end{myanswer}
\end{mychallenge}

\begin{mychallenge}[Typen von Rekurrenz und Transienz]
    Ordnen Sie für jeden Zustand der folgenden Markov-Ketten die Begriffe \emph{positiv rekurrent}, \emph{null-rekurrent} oder \emph{transient} zu. Begründen Sie Ihre Antwort jeweils kurz.

    \begin{enumerate}
        \item \textbf{Kette A:} \\
              $P = \begin{pmatrix}
                      0{,}5 & 0{,}5 \\
                      0{,}5 & 0{,}5 \\
                  \end{pmatrix}$

        \item \textbf{Kette B:} \\
              $P = \begin{pmatrix}
                      1     & 0     \\
                      0{,}3 & 0{,}7 \\
                  \end{pmatrix}$

        \item \textbf{Kette C:} (Zustände $1,2,3,\ldots$; von $i$ geht es immer mit Wahrscheinlichkeit $1$ zu $i+1$)

        \item \textbf{Kette D:} (Zustände $1,2,3,\ldots$; von $i$ geht es mit Wahrscheinlichkeit $1/2$ zu $i+1$ und mit $1/2$ zu $i-1$, wobei von $1$ aus nur nach $2$ möglich ist)
    \end{enumerate}

    \begin{myanswer}
        \begin{enumerate}
            \item \textbf{Kette A:} Beide Zustände sind \emph{positiv rekurrent}, da die Rückkehrzeiten endlich sind und die Kette endlich und irreduzibel ist.
            \item \textbf{Kette B:} Zustand 1 ist \emph{positiv rekurrent} (absorbierend), Zustand 2 ist \emph{transient}, da man von 2 mit positiver Wahrscheinlichkeit nach 1 gelangt und dann nie zurückkehrt.
            \item \textbf{Kette C:} Alle Zustände sind \emph{transient}, da man nie in einen Zustand zurückkehren kann (es geht immer nur vorwärts).
            \item \textbf{Kette D:} Alle Zustände sind \emph{null-rekurrent}, da man zwar mit Wahrscheinlichkeit 1 zurückkehrt, aber die erwartete Rückkehrzeit unendlich ist.
        \end{enumerate}
    \end{myanswer}
\end{mychallenge}

\begin{mychallenge}[Ergodizität]
    Welche Bedingungen müssen erfüllt sein, damit eine Markov-Kette ergodisch ist?

    \begin{myanswer}
        Alle Zustände müssen positiv rekurrent und aperiodisch sein.
    \end{myanswer}
\end{mychallenge}